\section{Inteligência Artificial Generativa e LLMs}

\textbf{OBJETIVO GERAL:}

Desenvolver competências para compreender, utilizar, adaptar e integrar modelos de Inteligência Artificial Generativa e Grandes Modelos de Linguagem (LLMs), aplicando técnicas modernas de engenharia de prompts, engenharia de contexto, recuperação aumentada por geração (RAG) e agentes inteligentes na construção de aplicações corporativas.

\textbf{CARGA HORÁRIA:} 180 horas

\textbf{FUNÇÃO:}

Desenvolver sistemas computacionais baseados em Inteligência Artificial, Análise Preditiva e Ecossistemas IoT, atendendo às normas de qualidade corporativas, robustez metodológica, governança de dados e arquitetura segura.

\textbf{SUBFUNÇÕES:}

\begin{itemize}
\item Compreender os fundamentos da Inteligência Artificial Generativa.
\item Desenvolver aplicações utilizando Grandes Modelos de Linguagem.
\item Projetar fluxos utilizando Engenharia de Prompt e Engenharia de Contexto.
\item Integrar modelos generativos em aplicações computacionais.
\item Desenvolver agentes inteligentes utilizando arquiteturas modernas.
\end{itemize}

\textbf{PADRÕES DE DESEMPENHO:}

\begin{itemize}
\item Identificar aplicações de IA Generativa.
\item Selecionar modelos adequados para diferentes problemas.
\item Desenvolver prompts estruturados.
\item Construir fluxos utilizando RAG.
\item Integrar bancos vetoriais às aplicações.
\item Desenvolver agentes inteligentes.
\item Avaliar respostas produzidas pelos modelos.
\item Aplicar princípios de IA Responsável.
\end{itemize}

\textbf{CAPACIDADES TÉCNICAS:}

\textbf{Fundamentos da IA Generativa}

\begin{itemize}
\item Diferenciar modelos discriminativos e generativos.
\item Identificar arquiteturas modernas de IA Generativa.
\item Reconhecer aplicações dos Grandes Modelos de Linguagem.
\item Avaliar limitações dos modelos generativos.
\end{itemize}

\textbf{Engenharia de Prompt}

\begin{itemize}
\item Elaborar prompts estruturados.
\item Aplicar técnicas de Zero-Shot, One-Shot e Few-Shot Prompting.
\item Desenvolver cadeias de raciocínio (Chain of Thought).
\item Refinar instruções para melhoria das respostas.
\end{itemize}

\textbf{Engenharia de Contexto}

\begin{itemize}
\item Organizar contexto para modelos generativos.
\item Gerenciar memória de conversação.
\item Estruturar conhecimento para recuperação de informações.
\item Integrar diferentes fontes de dados.
\end{itemize}

\textbf{RAG e Bancos Vetoriais}

\begin{itemize}
\item Construir pipelines de Retrieval-Augmented Generation.
\item Utilizar embeddings.
\item Integrar bancos vetoriais.
\item Avaliar qualidade da recuperação de informações.
\end{itemize}

\textbf{Agentes Inteligentes}

\begin{itemize}
\item Desenvolver agentes baseados em LLMs.
\item Integrar ferramentas externas.
\item Orquestrar fluxos multiagentes.
\item Automatizar tarefas inteligentes.
\end{itemize}

\textbf{CONHECIMENTOS:}

\textbf{Fundamentos da IA Generativa}

\begin{itemize}
\item Inteligência Artificial Generativa
\item Grandes Modelos de Linguagem (LLMs)
\item Modelos Foundation
\item Tokens
\item Janela de Contexto
\item Temperatura
\end{itemize}

\textbf{Arquiteturas}

\begin{itemize}
\item Transformers
\item Attention
\item Self-Attention
\item Decoder
\item Encoder
\item Embeddings
\end{itemize}

\textbf{Modelos}

\begin{itemize}
\item Llama
\item Mistral
\item Gemma
\item Qwen
\item DeepSeek
\item GPT (conceitos)
\end{itemize}

\textbf{Engenharia de Prompt}

\begin{itemize}
\item Zero-Shot Prompting
\item One-Shot Prompting
\item Few-Shot Prompting
\item Chain of Thought
\item Role Prompting
\item Prompt Chaining
\end{itemize}

\textbf{Engenharia de Contexto}

\begin{itemize}
\item Context Window
\item Context Engineering
\item Memória
\item Recuperação de Conhecimento
\item Estruturação de Contexto
\end{itemize}

\textbf{RAG}

\begin{itemize}
\item Retrieval-Augmented Generation
\item Chunking
\item Embeddings
\item Similaridade Vetorial
\item Re-ranking
\end{itemize}

\textbf{Bancos Vetoriais}

\begin{itemize}
\item ChromaDB
\item FAISS
\item Milvus
\item Pinecone (conceitos)
\end{itemize}

\textbf{Frameworks}

\begin{itemize}
\item LangChain
\item LangGraph
\item CrewAI
\item LlamaIndex
\item MCP
\end{itemize}

\textbf{Agentes Inteligentes}

\begin{itemize}
\item Agentes Baseados em LLM
\item Ferramentas (Tools)
\item Memória
\item Planejamento
\item Multiagentes
\end{itemize}

\textbf{Ética e Governança}

\begin{itemize}
\item IA Responsável
\item Hallucinations
\item Segurança em LLMs
\item Prompt Injection
\item Privacidade
\item LGPD
\end{itemize}

\textbf{CAPACIDADES SOCIOEMOCIONAIS:}

\begin{itemize}
\item Demonstrar criatividade na construção de soluções baseadas em IA Generativa.
\item Avaliar criticamente respostas produzidas por modelos generativos.
\item Formular estratégias para melhoria contínua dos prompts.
\item Trabalhar colaborativamente no desenvolvimento de aplicações inteligentes.
\item Produzir documentação técnica dos fluxos desenvolvidos.
\item Compartilhar conhecimento sobre arquiteturas generativas.
\item Demonstrar responsabilidade na utilização de modelos de IA.
\item Atuar observando princípios éticos, legais e de IA Responsável.
\end{itemize}

\textbf{AMBIENTES PEDAGÓGICOS:}

\begin{itemize}
\item Laboratório de Inteligência Artificial.
\item Laboratório de Desenvolvimento de Software.
\item Ambiente Virtual de Aprendizagem.
\item Ambiente Cloud para execução de modelos de IA Generativa.
\end{itemize}

\textbf{EQUIPAMENTOS E RECURSOS:}

\textbf{Hardware}

\begin{itemize}
\item Computadores com acesso à Internet.
\item Estações de trabalho com acesso a servidores de IA.
\item Projetor multimídia.
\end{itemize}

\textbf{Software}

\begin{itemize}
\item Python
\item Ollama
\item Open WebUI
\item LangChain
\item LangGraph
\item CrewAI
\item LlamaIndex
\item ChromaDB
\item Docker
\item VS Code
\item Jupyter Notebook
\item Git
\end{itemize}
